\section{Introduction}
\label{sec:introduction}

Diatoms fix roughly a fifth of global photosynthetic carbon, and in coastal
systems the rate at which they do so is usually set by dissolved inorganic
nitrogen rather than by light or phosphate. When nitrogen runs short a diatom
cell does not simply slow down. It reorganises: high affinity transporters are
placed in the plasma membrane, internal nitrogen is scavenged from pigment and
ribosomal protein pools, and carbon that would have gone into protein is
redirected into lipid and carbohydrate storage \citep{hockin2012response}.
The reorganisation is measurable in the transcriptome within hours, and the
order in which the modules respond tells us which of them are regulatory and
which are consequences.

\emph{Thalassiosira pseudonana} is the standard laboratory diatom for this
question because its genome is closed, well annotated, and small enough that a
modest sequencing budget gives saturating coverage \citep{armbrust2004genome}.
Strain \strain{} is a single-cell isolate maintained in the Coastal Genomics
Unit culture collection since 2019 and has been used in this practical series
for four cohorts.

This report has two aims. The first is analytical: to quantify which
transcripts change abundance between nitrogen-replete and nitrogen-limited
cultures, and to place the significant ones into functional groups. The second
is methodological: to document the sequencing quality control in enough detail
that a reader can judge whether the differential expression calls rest on sound
libraries. The second aim is not incidental. Two of our six libraries carried a
duplication rate above the level at which the teaching protocol normally
prescribes resequencing, and Section~\ref{sec:results} sets out the evidence on
which we decided to retain them.

We report three specific findings. Induction of the nitrate transporter family
is large, coherent across replicates, and detectable at every depth we
subsampled to. Repression of the photosystem II antenna is comparably large but
concentrated in a smaller number of transcripts. The urea cycle response is
weaker than either and would not have reached significance at the depth the
teaching protocol recommends, which is a useful negative result for anyone
planning a follow-up experiment.
